\clearpage
\begin{trapbox}
\textbf{误区一：把任意给定点当作切点}
\begin{center}
\textbf{$Q$ 不在曲线上，误套所得直线甚至不经过 $Q$}\par\medskip
\begin{tikzpicture}[x=.95cm,y=.95cm,>=Stealth,font=\small]
\draw[->,gray] (-4.7,0)--(5.3,0) node[right] {$x$};
\draw[->,gray] (0,-3.6)--(0,3.7) node[above] {$y$};
\draw[blue,very thick,domain=-1.43:1.43,samples=90] plot ({2*cosh(\x)},{sqrt(3)*sinh(\x)});
\draw[blue,very thick,domain=-1.43:1.43,samples=90] plot ({-2*cosh(\x)},{sqrt(3)*sinh(\x)});
\draw[red!75!black,dashed,thick] (4,-3.5)--(4,3.5);
\fill (1,0) circle (2pt);
\node at (1,1.4) {$Q(1,0)$};
\fill (4,3) circle (2pt) node[left=5pt] {$(4,3)$};
\fill (4,-3) circle (2pt) node[left=5pt] {$(4,-3)$};
\node[red!75!black,above] at (3.6,3.8) {误套公式：$x=4$};
\end{tikzpicture}

\end{center}
\textbf{错误做法：}对 $x^2/4-y^2/3=1$，将 $Q(1,0)$ 代入切点式，
得到 $x=4$ 并称它为 $Q$ 处的切线。

\textbf{为什么错：}$1^2/4-0^2/3=1/4\ne1$，所以 $Q$ 不在曲线上。
图中红色虚线与曲线交于 $(4,3)$、$(4,-3)$，是割线，也不经过 $Q$。

\textbf{正确替代：}先验点；若题目给的是切线经过的点，应另设切点，
像例 3 那样同时列出“切点在曲线上”与“切线经过给定点”两个条件。
\end{trapbox}
\clearpage
\begin{trapbox}
\textbf{误区二：记住外形，却颠倒了正负项}
\begin{center}
\textbf{同一双曲线 $y^2/4-x^2/3=1$，同一切点 $P(3,4)$}\par\medskip
\begin{tikzpicture}[x=.68cm,y=.68cm,>=Stealth,font=\small]
\foreach \shift/\kind in {0/0,9/1}{
\begin{scope}[shift={(\shift,0)}]
\draw[->,gray] (-3.7,0)--(3.9,0) node[right] {$x$};
\draw[->,gray] (0,-4.8)--(0,5.2) node[above] {$y$};
\draw[blue,very thick,domain=-1.5:1.5,samples=90] plot ({sqrt(3)*sinh(\x)},{2*cosh(\x)});
\draw[blue,very thick,domain=-1.5:1.5,samples=90] plot ({sqrt(3)*sinh(\x)},{-2*cosh(\x)});
\fill (3,4) circle (2pt) node[above left=2pt] {$P(3,4)$};
\ifnum\kind=0
\draw[red!75!black,dashed,thick] (-3.5,-2.71875)--(3.5,1.21875);
\node[align=center] at (0,-6) {错误：颠倒正负项\\$\frac{3x}{4}-\frac{4y}{3}=1$};
\else
\draw[orange!85!black,thick] (-3.5,-2.5)--(3.7,4.7);
\node[align=center] at (0,-6) {正确：保留正负项\\$y-x=1$};
\fi
\end{scope}}
\end{tikzpicture}

\end{center}
\textbf{错误做法：}写成 $3x/4-4y/3=1$，把正负项和分母机械套反。
代入 $P$，左边为 $-37/12\ne1$，连切点都不经过。

\textbf{正确替代：}原式的 $y^2/4$ 为正项，故替换为 $4y/4$；
$x^2/3$ 为负项，故替换为 $3x/3$，得到 $y-x=1$。

\textbf{误区三：漏掉竖直切线}

错误地给 $k=b^2x_0/(a^2y_0)$ 代入 $y_0=0$，会造成除以零。
应回到切点式，得到 $x=\pm a$；顶点图已经显示这两条竖直切线。

\medskip
\textbf{误区四：把“仅有一个公共点”当作相切的充分条件}

对 $x^2/4-y^2/3=1$，直线 $y=(\sqrt3/2)x+1$
与双曲线联立时，$x^2$ 项恰好抵消，只剩一次方程，
因此只有一个公共点 $(-4/\sqrt3,-1)$。
但该点处真正的切线斜率为 $\sqrt3$，与所给直线的斜率 $\sqrt3/2$ 不同，
所以它\textbf{不是切线}。
\textbf{正确做法：}判别式 $\Delta=0$ 只适用于联立后\textbf{二次项系数不为零}的情形；
若二次项消失，不能把“一次方程只有一个解”当作重根相切。
\end{trapbox}
